\documentclass[10pt,a4paper]{amsart}
\usepackage[margin=19mm]{geometry}
\usepackage{amsmath,amssymb,mathtools}
\usepackage[scr=rsfs,cal=euler]{mathalfa}
\usepackage{xfrac}
\usepackage[unicode,hidelinks,bookmarks=false]{hyperref}
\newcommand{\ie}{i.e.\ }
\newcommand{\cf}{cf.\ }
\newcommand{\resp}{respectively\ }
\newcommand{\Subsection}{\subsection}
\providecommand{\nobreakdash}{\nobreak\hbox{-}}
\setlength{\emergencystretch}{2em}
\multlinegap0pt
\usepackage{amssymb}
\usepackage{tikz}
\usepackage{float}
\newtheorem{theorem}{Theorem}
\title{Nijenhuis operators on 2D pre-Lie algebras and 3D associative algebras}
\author{Xiaoguang Zou, Xiang Gao, Chuangchuang Kang, Jiafeng L\"u}
\date{Source: arXiv:2308.12121v3 (CC BY 4.0)}
\begin{document}
\maketitle

\noindent\emph{Source and licence.} Nijenhuis operators on 2D pre-Lie algebras and 3D associative algebras. Version arXiv:2308.12121v3 (CC BY 4.0), distributed by arXiv under the Creative Commons Attribution 4.0 International licence, http://creativecommons.org/licenses/by/4.0/, as recorded in the arXiv licence metadata for this exact version and shown on the abstract page https://arxiv.org/abs/2308.12121v3. Full licence notice: \texttt{licenses/arxiv-2308.12121-CC-BY-4.0.txt}.\par\smallskip

\noindent\textit{Editorial scope.} Counted unit: the article's theorem c (source0.tex lines 709--740). The article's complete original proof of it is delivered with it (source0.tex lines 741--810, one \texttt{proof} environment of 70 lines). No mathematics is written, repaired or re-derived: the statement and the proof are the article's own lines, in the article's own notation, and no retained line is substituted or re-flowed.  Retained uncounted as the source-local context the proof needs: lines 663--665.  Nothing was omitted from any retained span, so the package carries no omission record.  No retained line contains a citation, so the wrapper carries no bibliography and no bibliography entry.  No retained line contains a figure environment, an image-inclusion command or drawing code, so no image is delivered and the package declares no asset dependency.  Editorial formatting only, in addition: an AMS \texttt{amsart} preamble that restates the article's own declarations and macros used by the retained lines, namely the environments \texttt{theorem} on separate counters, together with the standard packages those lines need.  The retained environments print in this document as Theorem 1 in order of appearance, whereas the article prints the same items with the numbering of its own full document.  The complete native source of the article as distributed in the arXiv e-print for this exact version is bundled unchanged as \texttt{sources/145998/source0.tex}.
\begin{equation}\label{eq:Nii-mulit}
  N\left(e_{i}\right)=\sum_{j=1}^{n} n_{i j} e_{j}, \quad ~n_{i j}\in\mathbb{C},
\end{equation}

\begin{theorem}\label{Ni ope on 2 c}
Let $N:A_i\rightarrow A_i,1\leq i\leq 5$, be a linear map on the pre-Lie algebra $(A_i, \cdot)$ defined by \eqref{eq:Nii-mulit}. Then the Nijenhuis operators on 2-dimensional commutative pre-Lie algebras are as follows:

$(1)$~The Nijenhuis operators on the pre-Lie algebra~$(A_{1},\cdot)$~are:
\begin{itemize}
   \item[] $N_{A_{1}}^1(e_1)= n_{11}e_1,~~ N_{A_{1}}^1(e_2)= n_{21}e_1+(n_{11}+n_{21})e_2,~n_{21}\neq0.$
   \item[] $N_{A_{1}}^2(e_1)= n_{11}e_1,~~ N_{A_{1}}^2(e_2)= n_{22}e_2.$
   \item[]  $N_{A_{1}}^3(e_1)=n_{11}e_1+ n_{12}e_2 ,~~ N_{A_{1}}^3(e_2)= (n_{11}-n_{12})e_2,~n_{12}\neq0.$
   \end{itemize}

$(2)$~The Nijenhuis operators on the pre-Lie algebra $(A_{2},\cdot)$~are:
\begin{itemize}
   \item[] $N_{A_{2}}^1(e_1)= n_{11}e_1,~~ N_{A_{2}}^1(e_2)= n_{22}e_2.$
   \item[]  $N_{A_{2}}^2(e_1)=n_{11}e_1+ n_{12}e_2 ,~~ N_{A_{2}}^2(e_2)= n_{11}e_2,~n_{12}\neq0.$
   \end{itemize}

$(3)$~The Nijenhuis operators on the pre-Lie algebra~$(A_{3},\cdot)$~are:
\begin{itemize}
   \item[] $N_{A_{3}}^1(e_1)= n_{11}e_1,~~ N_{A_{3}}^1(e_2)= n_{22}e_2.$
   \item[]  $N_{A_{3}}^2(e_1)=n_{11}e_1+ n_{12}e_2 ,~~ N_{A_{3}}^2(e_2)= n_{11}e_2,~n_{12}\neq0.$
   \end{itemize}

$(4)$~The Nijenhuis operators on the pre-Lie algebra~$(A_{4},\cdot)$~are:
\begin{itemize}
   \item[] $N_{A_{4}}^1(e_1)= n_{11}e_1+n_{12}e_2,~~ N_{A_{4}}^1(e_2)= n_{21}e_1+n_{22}e_2.$
   \end{itemize}

$(5)$~The Nijenhuis operators on the pre-Lie algebra~$(A_{5},\cdot)$~are:
\begin{itemize}
   \item[]  $N_{A_{5}}^1(e_1)=n_{11}e_1+ n_{12}e_2 ,~~ N_{A_{5}}^1(e_2)= n_{11}e_2.$
   \end{itemize}
\end{theorem}

   \begin{proof}
Let $\left\{e_{1},  e_{2}\right\}$ be a basis of 2-dimensional  commutative pre-Lie algebras~$(A_i,\cdot)$.~Set
  \begin{equation*}\label{eq:An operator on a basis}
    N(e_{1})=n_{11}e_1+n_{12}e_2,~N(e_{2})=n_{21}e_1+n_{22}e_2.
  \end{equation*}
Since $N$ is a Nijenhuis operator,~we have
 \begin{eqnarray}
    N(e_{1}) \cdot N(e_{1})&=&N(N(e_{1}) \cdot  e_{1}+e_{1} \cdot N(e_{1})-N(e_{1} \cdot  e_{1})),\label{eq:base e1e1}\\
     N(e_{1}) \cdot N(e_{2})&=&N(N(e_{1}) \cdot  e_{2}+e_{1} \cdot N(e_{2})-N(e_{1} \cdot  e_{2})),\label{eq:base e1e2}\\
      N(e_{2}) \cdot N(e_{1})&=&N(N(e_{2}) \cdot  e_{1}+e_{2} \cdot N(e_{1})-N(e_{2} \cdot  e_{1})),\label{eq:base e2e1}\\
       N(e_{2}) \cdot N(e_{2})&=&N(N(e_{2}) \cdot  e_{2}+e_{2} \cdot N(e_{2})-N(e_{2} \cdot  e_{2})).\label{eq:base e2e2}
 \end{eqnarray}
Calculating~\eqref{eq:base e1e1}~for the pre-Lie algebra~$(A_{1},\cdot)$, we have:
\begin{equation*}
% \nonumber % Remove numbering (before each equation)
   N(e_{1}) \cdot N(e_{1})
   %&= (n_{11}e_1+n_{12}e_2)(n_{11}e_1+n_{12}e_2)\nonumber  \\
   = {n_{11}}^2e_1+{n_{12}}^2e_2,\label{eq11}
\end{equation*}
and
\begin{equation*}
% \nonumber % Remove numbering (before each equation)
  N(N(e_{1}) \cdot  e_{1}+e_{1} \cdot N(e_{1})-N(e_{1} \cdot  e_{1}))
  % &=& N((n_{11}e_1+n_{12}e_2) \cdot  e_{1}+e_{1}\cdot(n_{11}e_1+n_{12}e_2)-N(e_1)) \nonumber\\
%  &=& N(2n_{11}e_1-n_{11}e_1-n_{12}e_2)\nonumber \\
%  &=& N(n_{11}e_1-n_{12}e_2)\nonumber\\
%  &=& {n_{11}}^2e_1+n_{11}n_{12}e_2-n_{12}n_{21}e_1-n_{12}n_{22}e_2 \nonumber \\
  = (n_{11}^2-n_{12}n_{21})e_1+(n_{11}n_{12}-n_{12}n_{22})e_2.\label{eq12}
\end{equation*}
Then
\begin{equation}\label{A11}
  n_{{11}}^2=n_{11}^2-n_{12}n_{21},~{n_{12}}^2=n_{11}n_{12}-n_{12}n_{22}.
\end{equation}
Similarly,~by (\ref{eq:base e1e2}) $\sim$ (\ref{eq:base e2e2}), we have:
\begin{eqnarray}
% \nonumber % Remove numbering (before each equation)
  n_{11}n_{21}=n_{11}n_{21}+n_{12}n_{21},~n_{12}n_{22}=n_{12}n_{21}+n_{12}n_{22}\label{A12}. \\
  {n_{21}}^2 = n_{21}n_{22}-n_{11}n_{21},~{n_{22}}^2={n_{22}}^2-n_{12}n_{21}\label{A13}.
\end{eqnarray}
Simplifying~\eqref{A11}$\sim$\eqref{A13}, we obtain the following three equations:
\begin{equation}\label{A14}
\left\{\begin{array}{l}
n_{12}n_{21}=0,\\{n_{12}}^2=n_{11}n_{12}-n_{12}n_{22},\\
{n_{21}}^2=n_{21}n_{22}-n_{11}n_{21}.
\end{array}\right.
\end{equation}

To solve the quadratic equations, we distinguish the two cases depending on whether
or not $n_{12} = 0$.

$\mathbf{Case~1}$: If $n_{12}=0$,~then~\eqref{A14} implies~$n_{21}(n_{21}-n_{22}+n_{11})=0$.~There are two subcases:~$n_{21}=0$ and $n_{21}\neq0$. If $n_{21}\neq0$,~we have~$n_{11}+n_{21}=n_{22}$.~Then we get ~$N_{A_{1}}^1$.
%:
%\begin{itemize}
%   \item[] $N_{A_{1}}^1(e_1)= n_{11}e_1,~~ N_{A_{1}}^1(e_2)= n_{21}e_1+(n_{11}+n_{21})e_2,~n_{21}\neq0.$
%   \end{itemize}
If $n_{21}=0$,~then we get $N_{A_{1}}^2$.
%:
%\begin{itemize}
%   \item[] $N_{A_{1}}^2(e_1)= n_{11}e_1,~~ N_{A_{1}}^1(e_2)= n_{22}e_2.$
%   \end{itemize}

$\mathbf{Case~2}$: If $n_{12}\neq0$, we have $n_{21}=0,~n_{22}=n_{11}-n_{12}$.~Then we get $N_{A_{1}}^3$.
%:
%\begin{itemize}
%   \item[]  $N_{A_{1}}^3(e_1)=n_{11}e_1+ n_{12}e_2 ,~~ N_{A_{1}}^3(e_2)= (n_{11}-n_{12})e_2,~n_{12}\neq0.$
%   \end{itemize}
Therefore, we obtain all Nijenhuis operators on the pre-Lie algebra $(A_1,\cdot)$.

Similarly, we can obtain all Nijenhuis operators on other 2-dimensional  commutative pre-Lie algebras. This completes the proof.
\end{proof}

\end{document}
